\documentclass[conference]{IEEEtran}
\usepackage{graphicx}
\usepackage{booktabs}
\usepackage{amsmath}
\usepackage{hyperref}
\title{EcoTrack: An Intelligent Building Energy Forecasting, Anomaly Detection, and Diagnostic Copilot Platform}
\author{\IEEEauthorblockN{Author One, Author Two, Author Three, Author Four, Author Five}\\\IEEEauthorblockA{Institution Name\\contact@example.com}}
\begin{document}
\maketitle
\begin{abstract}
EcoTrack combines load forecasting, anomaly detection, and a grounded conversational assistant to support building energy operations. This paper presents the system architecture, experimental protocol, and evaluation results.
\end{abstract}
\begin{IEEEkeywords}building energy management, load forecasting, anomaly detection, explainable AI, copilot
\end{IEEEkeywords}
\section{Introduction}
Buildings require timely, interpretable detection of abnormal electricity use. EcoTrack closes the loop from telemetry to operator diagnosis.
\section{System Architecture}
Figure~\ref{fig:architecture} summarizes the platform. Export the provided SVG to PDF before Overleaf compilation.
\begin{figure*}[t]
  \centering
  \includegraphics[width=0.96\textwidth]{../architecture/ecotrack-system-architecture.pdf}
  \caption{End-to-end EcoTrack system architecture.}
  \label{fig:architecture}
\end{figure*}
\section{Methodology}
Describe the BDG2 preprocessing pipeline, XGBoost confidence interval estimation, Isolation Forest detector, and grounded Copilot workflow.
\section{Experiments and Results}
Report train/test splits, forecasting error, anomaly metrics, latency, and ablation studies.
\section{Conclusion}
Summarize operational impact, limitations, and future work.
\bibliographystyle{IEEEtran}
\bibliography{references}
\end{document}
